% TOP_PROBLEM: {"id": 3166, "title": "Barnette's Conjecture", "category_id": 3, "category": {"id": 3, "name": "graph_theory", "display_name": "Graph Theory", "description": "Problems involving graphs, networks, and their properties.", "slug": "graph-theory", "order_index": 3, "created_at": "2026-07-31T15:26:25.671Z"}, "status": "open", "problem_number": "OPG-385", "set_id": 12}
% FIRST_POSTED: 2026-10-01
% ENTRY_KIND: statement_only
% UnsolvedMath ID 3166; source catalogue code OPG-385
% !TeX program = lualatex
% Definition and sources only; this file is not an LLM solution attempt.
\documentclass[11pt,a4paper]{article}
\usepackage[margin=25mm]{geometry}
\usepackage{fontspec}
\setmainfont{lmroman10-regular.otf}[BoldFont=lmroman10-bold.otf,
 ItalicFont=lmroman10-italic.otf,BoldItalicFont=lmroman10-bolditalic.otf]
\usepackage{amsmath,amssymb,mathtools}
\usepackage{unicode-math}
\setmathfont{latinmodern-math.otf}
\AtBeginDocument{\let\setminus\smallsetminus}
\usepackage{xurl,hyperref}
\hypersetup{colorlinks=true,urlcolor=blue}
\setcounter{secnumdepth}{0}
\setlength{\parindent}{0pt}
\setlength{\parskip}{5pt}
\setlength{\emergencystretch}{3em}
\begin{document}
\section*{Barnette's Conjecture}\label{3166}
{\small UnsolvedMath ID 3166 · OPG-385 · Graph Theory · OpenGarden}
\subsection{Definitions and mathematical statement}
All graphs in this statement are finite, simple, and undirected.
A graph $G=(V,E)$ is cubic if every vertex has degree three,
and bipartite if $V$ admits a partition into two sets with
every edge joining different sets. It is planar if its
vertices and edges can be drawn in the plane without crossings
between edge interiors. It is $3$-connected if $|V|\geq4$ and
removing any set of at most two vertices leaves it connected.

A Hamilton cycle is a simple cycle containing every vertex
of the graph. Equivalently, for $n=|V|$, it is a cyclic
ordering $v_1,\ldots,v_n$ of all vertices, without repetition,
such that $v_iv_{i+1}\in E$ for each index interpreted modulo
$n$.

Barnette's conjecture asserts that every graph satisfying
all four properties---$3$-connected, cubic, planar, and
bipartite---has a Hamilton cycle. The cycle uses exactly two
of the three incident edges at every vertex. Consequently,
its unused edges form a perfect matching, namely a set of
edges incident with every vertex exactly once. An equivalent
target is therefore to find a perfect matching whose edge
complement is one connected cycle.

There is no designated planar embedding or prescribed set
of edges the cycle must use. A union of several disjoint
cycles covering all vertices is only a spanning $2$-regular
subgraph and does not fulfil the required connectedness.
\subsection{Short English statement}
Does every 3-connected cubic planar bipartite graph have a cycle visiting every vertex exactly once?
\subsection{Sources}
\begin{itemize}
\item[S1] ULAM.AI and the UnsolvedMath Contributors, supplied problems.json, record 3166 (OPG-385), OpenGarden collection. Catalogue identity and source transcription; CC BY 4.0.
\url{https://huggingface.co/datasets/ulamai/UnsolvedMath}.
\item[S2] Open Problem Garden, Barnette's Conjecture. Original problem and discussion; the mathematical formulation below is expanded from the supplied source record.
\url{http://www.openproblemgarden.org/op/barnettes_conjecture}.
\end{itemize}
\end{document}
